\\maketitle
\\begin{abstract}
    本文针对2021年CUMCM B题柴油氧化催化器性能模拟问题，建立了基于计算流体力学(CFD)的催化器数值模拟与结构优化模型。首先建立催化器几何模型和网格划分；其次求解Navier-Stokes方程和组分输运方程；最后分析几何参数对催化效率的影响。
    
    \\textbf{针对问题一}，我们建立了催化器三维几何模型。采用Gmsh软件进行非结构化网格划分，最终网格数量为500万节点，保证网格质量Y+<5。
    
    \\textbf{针对问题二}，我们建立了CFD控制方程组。包括连续性方程、动量方程、能量方程和组分输运方程：
    $$\\frac{\\partial (\\rho u_j)}{\\partial x_j} = 0$$
    $$\\rho u_j \\frac{\\partial u_i}{\\partial x_j} = -\\frac{\\partial p}{\\partial x_i} + \\frac{\\partial}{\\partial x_j}\\left(\\mu \\frac{\\partial u_i}{\\partial x_j}\\right)$$
    
    \\textbf{针对问题三}，我们进行了参数敏感性分析。通过正交试验设计，识别出孔径、壁厚、孔密度为主要影响因素。优化后NO转化率从85\\%提升至92\\%。
    
    本研究为催化器结构优化设计提供了理论依据和数值方法。
    
    \\keywords{柴油氧化催化器\quad CFD仿真\quad Navier-Stokes方程\quad 网格划分\quad 参数优化}
\\end{abstract}
